\documentclass[11pt]{article}
\usepackage{amsmath, amssymb}
\usepackage{geometry}
\usepackage{array}
\usepackage{booktabs}
\geometry{margin=0.7in}
\setlength{\parindent}{0pt}
\setlength{\parskip}{0.35em}
\renewcommand{\arraystretch}{1.15}

\newcommand{\cheatsheettitle}[2]{%
\begin{center}
{\LARGE #1}\\[0.4em]
{\large #2}
\end{center}
}


\begin{document}
\cheatsheettitle{Algebra II Trigonometry Cheatsheet}{Module 9}

\section*{Angle Basics}
\[
180^\circ=\pi\text{ radians},\qquad 360^\circ=2\pi\text{ radians}
\]
\[
\theta^\circ=\theta\cdot\frac{\pi}{180},\qquad
\theta\text{ radians}=\theta\cdot\frac{180}{\pi}^\circ
\]

\section*{Special Values}
\[
\sin\left(\frac{\pi}{6},\frac{\pi}{4},\frac{\pi}{3}\right)=
\left(\frac12,\frac{\sqrt2}{2},\frac{\sqrt3}{2}\right)
\]
\[
\cos\left(\frac{\pi}{6},\frac{\pi}{4},\frac{\pi}{3}\right)=
\left(\frac{\sqrt3}{2},\frac{\sqrt2}{2},\frac12\right)
\]
\[
\tan\left(\frac{\pi}{6},\frac{\pi}{4},\frac{\pi}{3}\right)=
\left(\frac{\sqrt3}{3},1,\sqrt3\right)
\]
Quadrant signs:
\[
\text{QI all},\qquad \text{QII sin},\qquad \text{QIII tan},\qquad \text{QIV cos}
\]

\section*{Core Identities}
\[
\sin^2x+\cos^2x=1
\]
\[
1+\tan^2x=\sec^2x
\]
\[
1+\cot^2x=\csc^2x
\]
\[
\tan x=\frac{\sin x}{\cos x},\qquad
\cot x=\frac{\cos x}{\sin x}
\]

\section*{Symmetry and Cofunctions}
\[
\sin(-x)=-\sin x,\qquad \cos(-x)=\cos x,\qquad \tan(-x)=-\tan x
\]
\[
\sin\left(\frac{\pi}{2}-x\right)=\cos x,\qquad
\cos\left(\frac{\pi}{2}-x\right)=\sin x
\]
\[
\tan\left(\frac{\pi}{2}-x\right)=\cot x
\]

\section*{Angle Formulas}
\[
\sin(\alpha+\beta)=\sin\alpha\cos\beta+\cos\alpha\sin\beta
\]
\[
\sin(\alpha-\beta)=\sin\alpha\cos\beta-\cos\alpha\sin\beta
\]
\[
\cos(\alpha+\beta)=\cos\alpha\cos\beta-\sin\alpha\sin\beta
\]
\[
\cos(\alpha-\beta)=\cos\alpha\cos\beta+\sin\alpha\sin\beta
\]
\[
\tan(\alpha+\beta)=\frac{\tan\alpha+\tan\beta}{1-\tan\alpha\tan\beta}
\]
\[
\tan(\alpha-\beta)=\frac{\tan\alpha-\tan\beta}{1+\tan\alpha\tan\beta}
\]

\section*{Double and Half Angles}
\[
\sin(2x)=2\sin x\cos x
\]
\[
\cos(2x)=\cos^2x-\sin^2x=2\cos^2x-1=1-2\sin^2x
\]
\[
\tan(2x)=\frac{2\tan x}{1-\tan^2x}
\]
\[
\sin\left(\frac x2\right)=\pm\sqrt{\frac{1-\cos x}{2}},\qquad
\cos\left(\frac x2\right)=\pm\sqrt{\frac{1+\cos x}{2}}
\]
The sign depends on the quadrant of $\frac x2$.

\section*{Graphing}
\[
y=a\sin(b(x-h))+k
\]
\[
y=a\cos(b(x-h))+k
\]
\[
\text{amplitude}=|a|,\qquad \text{period}=\frac{2\pi}{|b|}
\]
\[
\text{midline }y=k,\qquad \text{phase shift }h
\]
For tangent:
\[
y=a\tan(b(x-h))+k,\qquad \text{period}=\frac{\pi}{|b|}
\]

\section*{Inverse Trig}
\[
\sin^{-1}x\in\left[-\frac{\pi}{2},\frac{\pi}{2}\right]
\]
\[
\cos^{-1}x\in[0,\pi]
\]
\[
\tan^{-1}x\in\left(-\frac{\pi}{2},\frac{\pi}{2}\right)
\]

\section*{Solving Trig Equations}
\begin{itemize}
  \item Isolate the trig expression.
  \item Find reference angle.
  \item Use quadrant signs.
  \item List all solutions in the requested interval.
\end{itemize}

\end{document}
